\documentclass[11pt]{amsart}
\usepackage[T1]{fontenc}
\usepackage[utf8]{inputenc}
\usepackage{amsmath,amssymb,amsthm}
\usepackage{mathtools}
\usepackage[margin=1.1in]{geometry}

\providecommand{\C}{\mathbb{C}}
\providecommand{\R}{\mathbb{R}}
\providecommand{\Z}{\mathbb{Z}}
\providecommand{\cA}{\mathcal{A}}
\providecommand{\cE}{\mathcal{E}}
\providecommand{\cO}{\mathcal{O}}
\providecommand{\cC}{\mathcal{C}}
\providecommand{\Obs}{\mathrm{Obs}}
\providecommand{\dbar}{\bar{\partial}}
\providecommand{\g}{\mathfrak{g}}
\providecommand{\tr}{\mathrm{Tr}}
\providecommand{\gh}{\mathrm{gh}}
\providecommand{\BV}{\mathrm{BV}}
\providecommand{\hCS}{\mathrm{hCS}}
\providecommand{\kappach}{\kappa_{\mathrm{ch}}}
\providecommand{\kappacat}{\kappa_{\mathrm{cat}}}
\providecommand{\Map}{\mathrm{Map}}
\providecommand{\PTVV}{\mathrm{PTVV}}
\providecommand{\CPTVV}{\mathrm{CPTVV}}
\providecommand{\Crit}{\mathrm{Crit}}
\providecommand{\Pois}{\mathrm{Pois}}

\theoremstyle{plain}
\newtheorem{theorem}{Theorem}[section]
\newtheorem{proposition}[theorem]{Proposition}
\newtheorem{lemma}[theorem]{Lemma}
\newtheorem{corollary}[theorem]{Corollary}
\theoremstyle{definition}
\newtheorem{definition}[theorem]{Definition}
\newtheorem{construction}[theorem]{Construction}
\newtheorem{remark}[theorem]{Remark}

\title{10D holomorphic Chern--Simons on a CY$_5$: the $(+1)$-shifted
Poisson partner}
\author{Wave-16 U4/20}
\date{2026-04-22}

\begin{document}
\maketitle

Wave-15~L4 established the BV shift law $n = d-4$ across the three
negative-shift strata $d \in \{2,3,4\}$, terminating the table at
$d=4$ where the shift crosses zero and the cotangent bracket
degenerates into a pointed $E_0$-datum. The termination was premature.
Pantev--To\"en--Vaqui\'e--Vezzosi \cite{PTVV2013} construct the
canonical $n$-shifted symplectic form on $\Map(Y_d, BG)$ for
\emph{every} integer $n$; when $Y_d$ is a Calabi--Yau $d$-fold the
shift is $n = d-4$, and \emph{positive} values $n \geq 1$ are
admissible provided the pairing is read as a shifted Poisson rather
than shifted symplectic structure on observables. The $d=5$ entry,
shift $+1$, is the first of this new regime.

\section{10D hCS on a CY$_5$}

Fix $X$ a compact Calabi--Yau 5-fold with nowhere-vanishing
holomorphic 5-form $\Omega_5 \in H^{5,0}(X)$, and $\g$ reductive with
invariant trace $\tr$. The fields are the Dolbeault ghost tower
\[
\cA \in \cE \;=\; \Omega^{0,\bullet}(X,\g)[1],\qquad
\cA \;=\; c + A_{0,1} + A^*_{0,2} + \cdots + A^{\#}_{0,4} + c^*_{0,5},
\]
with Dolbeault degrees $0, 1, 2, 3, 4, 5$ and ghost numbers
$+1, 0, -1, -2, -3, -4$ under the shifted-coordinate convention of
L4. The classical action is
\[
S_\hCS^{(10)}(\cA) \;=\; \int_X \Omega_5 \wedge \tr\!\Bigl(
 \tfrac12\, \cA \wedge \dbar\cA \;+\; \tfrac13\, \cA \wedge
 [\cA, \cA]_{\mathrm{Lie}_5} \Bigr).
\]
The subscript $\mathrm{Lie}_5$ records only that the associative
trace is cyclically invariant on matrix triples; on a Lie algebra the
cubic reduces to the standard wedge-Lie cube. Variation gives the
partial flatness equation
\[
\mathrm{EL}\bigl(S_\hCS^{(10)}\bigr) \;=\; \dbar A_{0,1} +
\tfrac12 [A_{0,1}, A_{0,1}] \;\in\; \Omega^{0,1}(X, \g),
\]
whose zero locus is $\Map_{\mathrm{hol}}(X_{\mathrm{dR}}, BG)$ in
$(0,1)$-gauge. The ten real dimensions of $X$ furnish the
$10$D $\hCS$ nomenclature, parallel to the $6$D / $8$D / $4$D strata
of L4.

\section{$(+1)$-shifted Poisson structure}

\begin{theorem}[PTVV shift at $d=5$]
\label{thm:plus1-shift}
The BV pairing
\[
\omega_\BV^{(5)} \;=\; \int_X \Omega_5 \wedge \tr\!\bigl(
\delta\cA \wedge \delta\cA\bigr)
\]
is a $(+1)$-shifted pairing on $\cE$. Under the identification
$\Crit(S_\hCS^{(10)}) = \Map(X_{\mathrm{dR}}, BG)$, the pullback of
the PTVV canonical form of \cite{PTVV2013}, Theorem~0.1 / §2, along
the AKSZ transgression over the CY$_5$ fundamental class
$[X] \in H_{5,5}(X)$ is $(+1)$-shifted.
\end{theorem}

\begin{proof}[Sketch]
Serre duality on $X$ pairs $\Omega^{0,k}(X,\g)$ with
$\Omega^{0,5-k}(X,\g^\vee)$ via $\Omega^{0,k} \otimes \Omega^{0,5-k}
\to \Omega^{0,5} \xrightarrow{\int \Omega_5 \wedge} \C$. The six-layer
tower collapses into three Serre-dual pairs:
$(c, c^*_{0,5})$ at Dolbeault degrees $(0,5)$;
$(A_{0,1}, A^{\#}_{0,4})$ at $(1,4)$;
$(A^*_{0,2}, A^\dagger_{0,3})$ at $(2,3)$. Each pair sums to total
Dolbeault degree $5$ and contributes to a pairing of ghost-number
$\gh = +1$ under the L4 convention: at $d=5$, ghost $c$ sits at
$\gh = +1$ and its dual $c^*_{0,5}$ at $\gh = -6$, with
$\gh(\delta c) + \gh(\delta c^*) - (2 \cdot \text{form-shift}) = +1$.
The degree equals $d - 4 = +1$, the PTVV shift law.
\end{proof}

The positivity of the shift is the substantive obstruction. PTVV §1.2
defines $n$-shifted symplectic as a closed nondegenerate $2$-form of
cohomological degree $n$; at $n = +1$ the nondegeneracy pairs
$\mathbb{T}_{\cE}$ with $\mathbb{T}_{\cE}[+1]$, which is not the
self-pairing required of a genuine symplectic structure. The correct
quantisation target is:

\begin{corollary}[Poisson not symplectic at $d=5$]
\label{cor:poisson-not-symplectic}
$\omega_\BV^{(5)}$ defines a $(+1)$-shifted \emph{Poisson} structure
on $\Obs_{\hCS}^{(10)}$ in the sense of Calaque--Pantev--To\"en--
Vaqui\'e--Vezzosi \cite{CPTVV2017}, not a symplectic one. The Poisson
bracket $\{-,-\}$ raises cohomological degree by $1$.
\end{corollary}

\begin{proof}
CPTVV §3 (Theorem~3.2.2 of the J.~Topology version) establishes the
equivalence $\Pois^{n}(-, \C) \simeq \mathrm{Symp}^{n}(-, \C)$ only
for $n \leq 0$; at $n \geq 1$ the left-hand side is genuinely richer
and the pairing is Poisson but not symplectic. The bracket
$\{-,-\}$ of \emph{positive} cohomological degree is the hallmark of
CY$_d$ BV for $d \geq 5$, contrasting with standard BV at
$d \leq 4$ where the bracket has degree $\leq 0$.
\end{proof}

\section{$\Obs_{\hCS}^{(10)}$ is an $E_5$-algebra}

\begin{theorem}[$E_5$-observables]
\label{thm:E5-obs}
The factorisation algebra $\Obs_{\hCS}^{(10)}$ of classical
observables of $S_\hCS^{(10)}$ on $X$ carries an $E_5$-algebra
structure on its global sections, with the bracket of
Corollary~\ref{cor:poisson-not-symplectic} as Poisson-$E_5$ bracket
of degree $+1$.
\end{theorem}

\begin{proof}[Sketch]
Costello--Gwilliam Vol~II, §4.8 \cite{CG2021}, shows that on a
$d$-complex-dimensional holomorphic factorisation geometry each
holomorphic direction contributes an $E_1$-factor after passing to
$\dbar$-cohomology; Lurie additivity
$E_1^{\otimes d} = E_d$ then gives $E_5$ at $d=5$. The Poisson
bracket from $\omega_\BV^{(5)}$ lifts to a Poisson-$E_5$ structure
on the cohomology of $\Obs_{\hCS}^{(10)}$ via the CPTVV
$\mathrm{BD}_{n+1} = E_n$-Poisson recognition (their §3.5).
\end{proof}

\section{Complete shift table, $d \in \{2,3,4,5\}$}

Appending the $d=5$ row to the L4 table:
\[
\begin{array}{c|cccc}
d & \text{hCS dim} & \text{PTVV shift} & \text{Obs-structure}
 & \text{Pairing type}\\
\hline
2 & 4 & -2 & E_2 & \text{symplectic}\\
3 & 6 & -1 & E_1 & \text{symplectic}\\
4 & 8 & \phantom{-}0 & E_0 & \text{symplectic (deg.)}\\
5 & 10 & +1 & E_5 & \text{Poisson (non-sympl.)}
\end{array}
\]
The $E_n$-index and PTVV shift no longer sum to $d-2$; the jump at
$d=5$ reflects the transition from cotangent-type BV (negative shift)
to Poisson-type BV (positive shift). Costello--Gwilliam additivity
counts the $E_n$-factor by holomorphic directions alone
($n = d$ at $d=5$), while the PTVV shift continues to measure
the Serre-dual imbalance $d - 4$.

\section{$\kappach$ on CY$_5$ examples}

\begin{proposition}[CY$_5$ supertrace values]
\label{prop:kappa-ch-CY5}
For the classical observable algebra at a compact CY$_5$ fixed point,
\[
\kappach(\cA_X) \;=\; \sum_{q=0}^{5} (-1)^q h^{0,q}(X)
\;=\; \chi(\cO_X) \;=\; \kappacat(X).
\]
Evaluation on two canonical examples:
\begin{itemize}
\item $X = \mathrm{K3} \times \mathrm{K3} \times E$: Künneth
multiplicativity gives $\kappacat(\mathrm{K3} \times \mathrm{K3}
\times E) = \chi(\cO_{\mathrm{K3}})^2 \cdot \chi(\cO_E) = 2 \cdot 2
\cdot 0 = 0$.
\item $X = $ a simply-connected compact CY$_5$ of ``quintic bundle''
type (the total space of a quintic polynomial bundle over a base,
cut to $\dim_\C X = 5$), with $h^{0,1}(X) = \cdots = h^{0,4}(X) = 0$
and $h^{0,0}(X) = h^{0,5}(X) = 1$ by Serre duality plus simple
connectivity: $\kappach = 1 - 0 + 0 - 0 + 0 - 1 = 0$ and $\chi(\cO_X)
= 1 + (-1)^5 = 0$. Every simply-connected CY$_5$ has vanishing
supertrace on the total space.
\end{itemize}
\end{proposition}

\begin{proof}
The supertrace identification $\kappach = \sum_q (-1)^q h^{0,q}$ is
the compact-CY$_d$ case of the Hodge-supertrace theorem of Vol~III
(cf.\ $\texttt{chapters/examples/cy\_d\_kappa\_stratification.tex}$
Theorem on the CY-D stratum). For $\mathrm{K3} \times \mathrm{K3}
\times E$, $\chi(\cO_E) = 0$ forces the product to vanish; the
triple-product echoes the K3$\times$E cancellation at $d=3$. For
the quintic, odd-dimensional CY forces $h^{0,0} - h^{0,5} = 1 - 1 = 0$
with intermediate Hodge numbers all zero.
\end{proof}

The K3$\times$K3$\times E$ cancellation is the $d=5$ echo of the
K3$\times E$ cancellation at $d=3$: any CY$_d$ containing an elliptic
factor has $\kappach = 0$ by the multiplicative vanishing
$\chi(\cO_E) = 0$.

\section{Wave-15~M3 intimation and the $\Lambda_{32}$ embedding}

Wave-15~M3 \cite{WaveM3} intimated a $\Lambda_{32} \hookrightarrow
\Lambda_{33}$ lattice embedding parameterising CY$_5$ data at the
K3$\times$K3$\times E$ point: the product
Aut of $(\mathrm{K3} \times \mathrm{K3})$ acts on the 32-dimensional
Mukai-squared lattice, with the elliptic factor contributing the
33rd basis vector as the spectral parameter. The $(+1)$-shifted
Poisson structure of Theorem~\ref{thm:plus1-shift} is compatible with
this embedding: the positive shift reflects the extra holomorphic
direction (the elliptic $E$) that does \emph{not} enter the
Mukai pairing on the two K3 factors, and the Poisson bracket of
Corollary~\ref{cor:poisson-not-symplectic} acts along this single
extra direction. The 10D hCS model therefore supplies the BV
underpinning of the M3 arithmetic structure, and the $(+1)$-shift
is the precise reason the M3 embedding is of codimension one rather
than being a genuine symplectic match.

At $d=5$ the CY-D dimensional stratification of $\Phi$ places
$\Phi_5(\mathrm{Coh}(X))$ in the $E_1$-chiral class (by the
$d \geq 3 \Rightarrow E_1$ law of Vol~III), while the 10D hCS
observables are $E_5$ on the BV side. The two statements are
compatible: $\Phi$ reads out the highest locally-constant direction
(one, along $\C$), whereas Costello--Gwilliam additivity counts all
five holomorphic disc-directions of $X$. The $E_5 \to E_1$ projection
is Lurie-Koszul forgetting four disc-directions.

\begin{thebibliography}{9}
\bibitem{PTVV2013} T.~Pantev, B.~To\"en, M.~Vaqui\'e, G.~Vezzosi,
  \emph{Shifted symplectic structures}, Publ.~IHES 117 (2013),
  271--328.
\bibitem{CPTVV2017} D.~Calaque, T.~Pantev, B.~To\"en, M.~Vaqui\'e,
  G.~Vezzosi, \emph{Shifted Poisson structures and deformation
  quantization}, J.~Topology 10 (2017), 483--584.
\bibitem{CG2021} K.~Costello, O.~Gwilliam, \emph{Factorization
  Algebras in Quantum Field Theory, Volume II}, Cambridge University
  Press (2021).
\bibitem{WaveM3} Wave-15 M3, \emph{$\Lambda_{32} \hookrightarrow
  \Lambda_{33}$ embedding for the K3$\times$K3$\times$E stratum},
  \texttt{notes/wave15\_m3\_Lambda32\_to\_33\_embedding.tex} (2026).
\end{thebibliography}

\end{document}
